% =====================================================================
%  CHAPTER 7: 儿童少年健康问题和健康促进策略 + 心理卫生问题
%  对应大纲第六章（心理问题）+ 往届笔记第八章（健康问题和健康促进策略）+ 往届笔记第十二章（心理卫生问题）
% =====================================================================
\chapter{儿童少年健康问题和健康促进策略}
\setcounter{chapter}{7}
\chapterintro{
\textbf{大纲要求：}
\begin{itemize}[leftmargin=*, itemsep=2pt]
    \item \textbf{掌握}（双横线）：心理卫生基本概念；心理健康标准和工作目标
    \item \textbf{熟悉}（单横线）：常见心理卫生问题
\end{itemize}
\vspace{4pt}
本章综合儿童少年健康问题和健康促进策略（往届笔记第八章）与儿童少年心理卫生问题（往届笔记第十二章，对应大纲第六章），分为两大部分：(A) 健康问题和健康促进策略，涵盖健康维度、主要健康问题、生命历程观、DOHaD理论及健康促进策略；(B) 心理卫生问题，涵盖心理卫生概述、ADHD、SLD、ASD等常见心理障碍。
}

% =====================================================================
%  PART A: 健康问题和健康促进策略
% =====================================================================
\section{衡量健康的维度和指标体系}
\label{sec:health-dimensions}

\subsection{健康的四个维度}

\begin{defbox}
\textbf{健康（Health）：}不仅仅是没有疾病或虚弱，而是一种身体、心理、社会适应和道德上的完好状态。
\end{defbox}

\begin{keybox}
\textbf{健康的四个维度及评价指标：}

\begin{enumerate}[leftmargin=*]
    \item \imp{生理健康}：体格发育、生理功能、运动能力、日常生活自理能力。
    \item \imp{心理健康}：心理/精神症状和行为表现、认知功能、智力和个性、主观幸福度。
    \item \imp{社会健康}。
    \item \imp{道德健康}。
\end{enumerate}
\end{keybox}

% ----- A.2 主要健康问题 -----
\section{儿童少年主要健康问题}
\label{sec:health-problems}

\begin{keybox}
\textbf{当前儿童少年主要健康问题：}

\begin{enumerate}[leftmargin=*]
    \item \imp{伤害成为儿童少年第一位死因}。
    \item \imp{心理行为问题日益突出}。
    \item \imp{肥胖和近视成为影响儿童少年健康的常见疾病}。
    \item \imp{慢性疾病低龄化日渐突显}。
    \item \imp{其他}：体能下降、传染病易感。
\end{enumerate}
\end{keybox}

\subsection{社会变革与中国儿童少年新健康问题}

\begin{defbox}
\textbf{（一）城市流动儿童：}各种传染性疾病发病风险增高、营养性疾病检出率高、心理压力大、健康危险行为多。

\textbf{（二）农村留守儿童：}心理健康问题、社会化和教育问题、安全问题、营养性疾病检出率高。
\end{defbox}

% ----- A.3 生命历程观 -----
\section{儿童少年健康的生命历程观}
\label{sec:life-course}

\subsection{健康的生命历程观}

\begin{keybox}
\textbf{生命历程（Life Course）：}生命过程的各个阶段组成了人的生命历程。

\textbf{生命历程观（Life Course Perspective）：}主张在多学科、多维度原则下，观察个体生活的物质和社会历史环境，考察各种社会经济、社会组织和制度、公共卫生及社区和家庭环境因素对个体身心健康的影响。
\end{keybox}

\subsection{人体功能水平的生命历程变化规律}

\begin{defbox}
\textbf{人体功能水平变化：}抛物线，逐渐上升 \(\to\) 持续一段高水平状态 \(\to\) 逐渐下降。

\textbf{生命历程的健康促进：}

\begin{enumerate}[leftmargin=*]
    \item \imp{生命早期}：包括围生期、婴幼儿期、青少年期，促进这一时期的生长发育，可使个体生理和心理功能水平尽可能到达高位。
    \item \imp{中年时期}：人体功能水平通常达到并维持高位，然后出现缓慢下降过程。因此这一时期应将人体功能水平尽可能维持在高位。
    \item \imp{老年时期}：人体功能水平快速下降，因此这一时期应尽可能减少其下降速度，让人能享有足够长的健康寿命，自主生活，减少残疾和失能困扰。
\end{enumerate}
\end{defbox}

\subsection{生命历程中慢性病危险度长期积累现象}

\begin{defbox}
慢性病发生与发展以及人体的生理功能改变是危险因素长期累积暴露的结果。危险因素包括基因、母亲孕期以及婴幼儿期的营养状况、家庭环境和社会关系的影响、个人的生活习惯和成年期的工作环境等。慢性非传染性疾病的累积危险度随年龄增长逐渐增多；危险作用可以发生在生命全程的所有阶段；虽然慢性病的积累危险度在成年后快速增高，但这是从生命早期阶段逐渐积累的结果。
\end{defbox}

% ----- A.4 健康危险因素的生命历程观 -----
\section{儿童少年健康危险因素的生命历程观}
\label{sec:risk-factors}

\subsection{母亲孕期健康与子女健康的关系}

\begin{keybox}
\textbf{健康和疾病的发育起源（Developmental Origins of Health and Disease, DOHaD）学说/"DOHaD理论"：}人体若在早期发育过程中（包括胎儿、婴幼儿时期）经历不利因素，比如营养不良或营养过剩，人体组织器官在结构和功能上将发生永久性或程序性改变，导致成年期一些慢性非传染性疾病（包括心血管疾病、糖尿病、肿瘤、哮喘、骨质疏松、代谢综合征、神经性疾病、精神障碍等）发生发展。

\textbf{成人疾病的胎源假说（Hypothesis of Fetal Origins of Adult Disease, FOAD）：}胎儿宫内不良反应使其自身代谢和器官组织结构发生适应性调节，若营养不良得不到及时纠正，这种适应性调节将导致重要机体组织和器官，包括血管、胰腺、肝脏和肺脏等，在代谢结构上永久性改变，进而演变为成人期疾病。
\end{keybox}

\subsection{儿童少年时期体格生长发育因素与成年期健康的关系}

\begin{defbox}
\textbf{影响成年期健康的儿童少年时期因素：}营养不良/过剩、儿童期肥胖、青春发动时相提前、幼年时期血压轨迹。
\end{defbox}

\subsection{形成于青少年时期的不健康心理行为问题}

\begin{defbox}
\textbf{青少年健康危险行为（Adolescent Health-Risk Behaviors）：}任何能给青少年健康和完好状态乃至成年期健康和生活质量造成直接或间接损害的行为。

\textbf{健康危险行为对人体健康损害作用的表现特点：}潜伏期长、特异性差、联合作用、广泛存在。
\end{defbox}

% ----- A.5 健康促进策略 -----
\section{儿童少年健康促进策略}
\label{sec:promotion-strategies}

\subsection{全生命周期保健策略}

\begin{keybox}
\textbf{全生命周期保健（Life Course Care）：}把生命体看作一个由出生、生长发育、成熟和衰老渐变的周期过程，通过把人生划分为胎儿期和婴幼儿期、儿童少年期、成年工作期和晚年期4个明确阶段，并针对这些不同年龄组人群在不同场所（家庭、托幼机构/学校、社区、工作场所）实施连续性预防服务措施，从而保证人生不同阶段有效获得有针对性的预防服务，减少不必要的重复和遗漏，既高效又节省地达到促进人群健康目的。这被认为是生命周期保健策略强调保证整个人群健康，促进健康老龄化的最佳途径。
\end{keybox}

\begin{defbox}
\textbf{全生命周期保健的特点：}

\begin{enumerate}[leftmargin=*]
    \item \imp{连续性}：生命形成前即孕前至整个生命的各个时期，生长发育的每一阶段，都是前一阶段的延续、后一阶段的开始，受前一阶段影响，同时又影响下一阶段。
    \item \imp{累积性}：自环境和个人行为因素在不同时点对人体健康产生长远影响。
    \item \imp{关键窗口}：特定关键/敏感期如胎儿期、幼儿期等，一些经历会影响一个人未来健康和发展，一些不良事件或暴露对个体产生的消极影响更强烈。
    \item \imp{意义}：在生命全程各个关键/敏感时期，明确或强化通过哪些保健和健康促进措施来改善整体健康。
\end{enumerate}
\end{defbox}

\subsection{生命初始1000天}

\begin{keybox}
\textbf{生命初始1000天（The First 1000 Days）：}从怀孕胚胎形成到出生后2岁的时期。建立在DOHaD理论上，强调这一时期健康的重要性。因为这一时期的营养，不仅关系到胎婴儿期体格发育和神经系统发育，还可能关系到成人后的健康。成年期心血管疾病、糖尿病、肥胖等都与胎儿、婴幼儿时期营养不良或过剩相关。该策略倡导人们重视生命最初1000天内健康生长发育，减少成年后慢性非传染病发生，促进并改善人类健康。
\end{keybox}

\subsection{基于健康社会决定因素理论的健康不平等改善策略}

\begin{defbox}
\textbf{健康不平等（Health Inequality）：}不同社会经济特征人群间的健康水平差异，已成为影响全球健康的核心问题。

\textbf{健康不公平（Health Inequity）：}从广义看健康不平等既包括健康状况差异，又包括获得良好健康状态的机会、条件的差异，即健康不公平。
\end{defbox}

\begin{defbox}
\textbf{决定儿童少年健康的社会环境因素：}

\begin{enumerate}[leftmargin=*]
    \item \imp{日常生活环境}。
    \item \imp{社会结构性因素}。
\end{enumerate}

\textbf{实现健康公平的行动原则：}

\begin{enumerate}[leftmargin=*]
    \item \imp{改善日常生活环境}。
    \item \imp{改善社会结构性因素}。
    \item \imp{提高公众认知}。
\end{enumerate}
\end{defbox}

\subsection{将健康融入所有政策的策略}

\begin{keybox}
\textbf{将健康融入所有政策（Health in All Policies, HiAP）：}旨在改善人群健康和健康公平的公共政策制定方法，系统考虑了公共政策可能带来的健康影响，寻求部门间合作，避免政策对公众健康造成不良影响。

2016年8月召开的全国卫生与健康大会上，\textbf{"将健康融入所有政策"}被确定为新时期我国卫生与健康工作的重要方针之一。
\end{keybox}

\newpage
